% !TEX root = userguide.tex

\chapter{Writing new source code for \tfem}
\pgst{tfemch}

\def\chaptertitle{Writing new source code for \tfem}
\def\headdate{14th December 2025}

In this chapter we will describe some tools and techniques for writing new
source code, including some info on how to write new elements and some special topics for writing main programs.

\section{Some tools for writing new code}
\subsection{The command \texttt{fmm}: make a makefile}
\label{sec:fmm}

The command \texttt{fmm} (fortran make makefile) is a simple shell script
to create a Makefile 
for all Fortran files
in the current directory to deal with module dependencies. Two modes of
operation are available:
\begin{quote}
       \texttt{fmm >Makefile}
\end{quote}
which produces a makefile for building main programs (when present) and
compiling the source files in the current directory, and
\begin{quote}
       \texttt{fmm -l >Makefile}
\end{quote}
which produces a makefile for producing a library (archive)
and thus no main programs are build.
Type \texttt{fmm -h} to see all the options available.

There are some restrictions on the use of \texttt{fmm}:
\begin{itemize}
   
   \item Several macros in the makefile are undefined. Although some are optional, they need to be defined 
    by the user. In \tfem\ the include file \texttt{Mdefs.mk} is used for that.
   \item Main programs must have an \texttt{end program} statement, otherwise 
      main programs are undetected and even might end up in a library.
   \item A module must be in a separate file or part of the main program file.
     The filename (without the extension) must be identical to the module name
     or the module name without the \texttt{\_m} extension. This
     convention is used throughout \tfem, i.e.\ all modules have a \texttt{\_m}
     extension and defined in separate files.
   \item Only object (\texttt{.o}) files produced by the Makefile are linked.
     So just adding an object file to the current directory will not work.
%     Note, that this is different from \sep, where all object files
%     in the current directory are linked into the executable.
   \item The main programs are linked to \emph{all} object (\texttt{.o})
     files produced by the Makefile.

\end{itemize}
A command \texttt{fmme} is available that removes the restriction
on modules (separate files with filename equal to the module name) and can be
used for a wider range of projects. However, it is several times slower than
\texttt{fmm} because of the searching for module definitions in the source
files.

The commands \texttt{fmm} and \texttt{fmme} are for development of source code
\emph{outside} the \tfem-tree. A special version of \texttt{fmm}, called
\texttt{mm}, is available for use \emph{inside} the \tfem-tree. It is not
available by default and can be created by making a symbolic link in the
\texttt{bin} subdirectory of the \tfem-tree:
\begin{quote}
   \texttt{ln -s fmm mm}
\end{quote}

\subsection{The command \texttt{showmodules}: show contents of modules}

The command \texttt{showmodules} will give for each module available a list of
the derived types and the subroutines it contains to the standard output.
Typing
\begin{quote}
       \texttt{showmodules >list}
\end{quote}
and editing the file \texttt{list} will give an easy access to the source using
the tags mechanism.

The command \texttt{showmodules} can also be used for individual source files
as follows
\begin{quote}
       \texttt{showmodules file(s)}
\end{quote}

\subsection{The \texttt{tags} file}

If all goes well and your source compiles/links, a file \texttt{tags} will be
produced for easy jumping around in the your own source and \tfem\ source. If
you do not have your own source files yet (or even an empty directory)
this will also work, so
\begin{quote}
   \texttt{getmdefs >Mdefs.mk}\\
   \texttt{fmm >Makefile}\\
   \texttt{make}
\end{quote}
will produce a \texttt{tags} file containing tags to the source of \tfem.

If your files do not compile/link no \texttt{tags} file is produced. Typing
\begin{quote}
   \texttt{make tags}
\end{quote}
will always make a new \texttt{tags} file.

\subsection{The command \texttt{tfgrep}: \texttt{grep} in source files}

The command \texttt{tfgrep} will do a `grep' on all \tfem\ source files. For example, type
\begin{quote}
   \texttt{tfgrep PATTERN}
\end{quote}
to search for the string \texttt{PATTERN} in the source files of
\texttt{src/core} and \texttt{src/addons}. Type
\begin{quote}
   \texttt{tfgrep -h}
\end{quote}
to show all the options available.

\subsection{The command \texttt{getsrc}: get source files}

The command \texttt{getsrc} will copy \tfem\ source files to the current
directory, for example
\begin{quote}
   \texttt{getsrc system.f90 vector.f90}
\end{quote}
will copy the two specified source files.

\subsection{The commands \texttt{tfdiff} and \texttt{tfvimdiff}: difference with \tfem\ files}
\label{sec:tfdiff}

The command \texttt{tfdiff} will show the difference between the
\tfem\ source files to files in the current
directory, for example
\begin{quote}
   \texttt{tfdiff system.f90 vector.f90}
\end{quote}
will show the difference of these files in the `unified diff' format.

The command \texttt{tfvimdiff} will show the difference between the
\tfem\ source files to files in the current
directory using the \texttt{vimdiff} utility, for example
\begin{quote}
   \texttt{tfvimdiff system.f90 vector.f90}
\end{quote}
will show the difference of these files.
Adding the option \texttt{-g}, i.e.\ \texttt{tfvimdiff -g},
will use the GUI-based \texttt{gvimdiff} instead of 
the text based \texttt{vimdiff}. Similarly, the option \texttt{-m} will use the
MacVim based \texttt{mvimdiff}.

For both \texttt{tfdiff} and \texttt{tfvimdiff} adding the option \texttt{-r} will 
reverse the comparison.

\section{Changing \tfem-files}

Sometimes you would like to change files in the \tfem-tree.
For example,
if you find a bug\footnote{Please report bugs to \texttt{m.a.hulsen@tue.nl}. If you know how to fix it, please send a patch with a
 fix (see Sec.~\ref{sec:patches}).}, want to change a subroutine of an element module or add a
new subroutine to it.
There are basically two ways to change or add files in \tfem:
\begin{itemize}
   \item Change or add the files in the \tfem-tree directly, optionally using a branch in git.
         This is only recommended for minor changes such as temporary bug fixes. Remember that in the \tfem\ tree \texttt{mm} should be used instead of \texttt{fmm} (see Sec.~\ref{sec:fmm}).
   \item Make a new module\footnote{Use a name for the module that does not yet exist in \tfem.} in your own working directory and `use' the
         \tfem-modules and add or change routines. This is typically used for
         research code that is not (yet) scheduled to be included in \tfem.
         For example, you want to add your own routines to the
         \texttt{stokes\_elements\_m} module (see
         Sec.~\ref{sec:stokesaddon}) and want also to change 
         the routine \texttt{stokes\_elem}. Now make a new module file
         containing:  
\begin{verbatim}
    module my_stokes_elements_m
      use stokes_elements_m
    contains              
       <your element routines including a new routine my_stokes_elem>       
    end module my_stokes_elements_m
\end{verbatim}
         Now you can use your own \texttt{my\_stokes\_elem} as a replacement for
         \texttt{stokes\_elem}. The latter is still available. If you want to keep the name
         \texttt{stokes\_elem} for your new routine, use
\begin{verbatim}
    module my_stokes_elements_m
      use stokes_elements_m, stokes_elem_tfem => stokes_elem
    contains              
      <your element routines including a new routine stokes_elem>   
    end module my_stokes_elements_m
\end{verbatim}
    where the routine \texttt{stokes\_elem} in \tfem\
         has been renamed to \texttt{stokes\_elem\_tfem} and you can use \texttt{stokes\_elem} yourself.
         
    If you now put the line 
\begin{verbatim}
      use my_stokes_elements_m
\end{verbatim}
    in your main program, the stokes routines in \tfem\ and your own
         routines will be available. 
    If your own project becomes too unwieldy with lots of files, you might want to have a look at Appendix~\ref{chap:projecttree} to make your own project tree with subdirectories.
 
\end{itemize}
\newpage
\section{Contributing to \tfem}

If you want to contribute code to \tfem, such as a fix for a bug or a new example or add-on, you can do that by 
either making a patch or by creating a branch and ask for a merge request on gitlab. 

\subsection{Making a patch}
\label{sec:patches}

You make a patch by cloning a fresh \tfem-tree and modify the tree to represent the changes you want to make. Go to the root of the \tfem-tree and type
\begin{quote}
   \texttt{git diff >mypatch.diff}
\end{quote}
or if you want to create a patch only for a subdirectory of the \tfem-tree type
\begin{quote}
   \texttt{git diff subdir >mypatch.diff}
\end{quote}
Compress the patch before sending it as an attachment by email:
\begin{quote}
   \texttt{gzip mypatch.diff}
\end{quote}
Please send the patch to \texttt{m.a.hulsen@tue.nl} with a short description of the changes.

\subsection{Creating a branch}
\label{sec:branch}

In order to contribute via a branch you need to be a member of the \texttt{tfem\_develop} project on the GitLab server \texttt{https://gitlab.tue.nl}.
Also, make sure you have configured your name and email correctly by using \texttt{git config}:
\begin{verbatim}
   git config --global user.name "John Doe"
   git config --global user.email johndoe@example.com
\end{verbatim}
To create a new branch, first pull a fresh \tfem-tree from \texttt{tfem\_develop} to your local computer.
Create a branch, edit the branch, make commits, push the branch to the server when you are finished and make a merge request on GitLab. Please follow the commit guidelines (in Chapter ``Distributed Git'') from the Pro Git book on the Git website \texttt{https://git-scm.com}. In particular, the commit message should have a single line with a maximum of around 50 characters. More detailed info can be written after an empty line following the first line. Each line in the detailed info should be wrapped at roughly 72 characters.

\newpage
\section{Programming new elements}

As already seen in the examples in Chap.~\ref{ch:examples} element routines are
treated by \tfem\ as if they are function routines that need to be written
outside the scope of the \tfem\ routines. The only thing the user has to do is
supply the name of the element routine via the \texttt{elemsub} or
\texttt{elemsub1} dummy heading parameter. The element routine must, of course,
conform to the \texttt{interface}
statement given in the heading definition of the \tfem\ routine that is called,
but otherwise the user is completely free in designing the elements.
Nevertheless, some guidelines and routines available for writing elements are
given here. The examples in the \tfem-tree should get you going after that.

\subsection{A general outline of an element routine}

A general outline of an element for \texttt{build\_system} that is properly
embedded in a module with the proper \texttt{use} statements, is for example

\begin{listing}{1}
module elements_m

  use tfem_elem_m

  implicit none

\end{listing}
\begin{quote}
   \emph{global definitions of this module}
\end{quote}
\begin{listingcont}

contains

  subroutine elemsub ( mesh, problem, elgrp, elem, matrix, vector, &
    first, last, coefficients, oldvectors, elemmat, elemvec )
    type(mesh_t), intent(in) :: mesh
    type(problem_t), intent(in) :: problem
    integer, intent(in) :: elgrp, elem
    logical, intent(in) :: matrix, vector, first, last
    type(coefficients_t), intent(in) :: coefficients
    type(oldvectors_t), intent(in) :: oldvectors
    real(dp), intent(out), dimension(:,:) :: elemmat
    real(dp), intent(out), dimension(:) :: elemvec

\end{listingcont}
\begin{quote}
   \quad\emph{Local definitions}
\end{quote}
\begin{listingcont}

    if ( first ) then
   
\end{listingcont}
\begin{quote}
   \qquad\parbox{10cm}{\emph{Allocate memory. Compute variables that are 
     constant for all elements in this group, for example shape functions.}}
\end{quote}
\begin{listingcont}

    end if
   
\end{listingcont}
\begin{quote}
   \quad\emph{Find coordinates and deformed shape of the element.}
\end{quote}
\begin{listingcont}

    if ( vector ) then
   
\end{listingcont}
\begin{quote}
   \qquad\emph{compute element right-hand side vector}
\end{quote}
\begin{listingcont}

    end if
   
    if ( matrix ) then
   
\end{listingcont}
\begin{quote}
   \qquad\emph{compute element matrix}
\end{quote}
\begin{listingcont}

    end if
   
\end{listingcont}
\begin{listingcont}

    if ( last ) then
   
\end{listingcont}
\begin{quote}
   \qquad\emph{deallocate memory for this group.}
\end{quote}
\begin{listingcont}

    end if
   
  end subroutine elemsub

end module elements_m
\end{listingcont}
\bigskip
Some notes:
\begin{description}
   \item[line 3] types and routines from other modules are `included'
   here and available everywhere in the module.
  \item[line 10-19] the heading can be copied from the \texttt{interface}
    definition of \texttt{elemsub} in the heading of \texttt{build\_system}.
    The subroutine name \texttt{elemsub} should preferable
    be changed to something else. If there is more than one routine you have
    to change it.
  \item[line 14] the integers \texttt{elgrp} (element group number)
   and \texttt{elem} (element number within this group) determine the 
   element\footnote{A new element interface \texttt{elemsub1} is 
   available that uses a structure to supply element information. This new
interface give more possibilities, such as building a system on objects.}.
  There is no global element number like in \sep.
  \item[15, 22 and 39] the logicals \texttt{first} and \texttt{last} are
   \texttt{.true.}\ when the current element is the first or the last element in
      this element group, respectively.
  \item[15, 28 and 33] the logicals \texttt{matrix} and \texttt{vector} are
    determined by the heading parameters \texttt{buildmatrix} and
    \texttt{buildvector}, respectively.
   \item[line 16] \texttt{coefficients} is a structure of type
    \texttt{coefficient\_t} for supplying parameters and coefficients to the
     element subroutine. The structure contains only data when the structure is 
     part of the heading of \texttt{build\_system}.
   \item[line 17] \texttt{oldvectors} is a structure of type
    \texttt{oldvectors\_t} for supplying already available vectors of type
     \texttt{sysvector\_t} and \texttt{vector\_t} to the
     element subroutine. The structure contains only data when the structure is 
     part of the heading of \texttt{build\_system}.
\end{description}

\subsection{Transferring data between element routines:
a vector defined per element} 
\label{sec:elvector}

Suppose we have to solve a diffusion problem where the diffusion coefficient
is not constant and depends on the position in a complicated way, possibly
depending on the solution of other problems. We can of course build a cathedral
and make a complicated diffusion element with all the bells and whistles needed
by the user. Another approach is to make a relatively simple diffusion element
with varying coefficients, where the latter needs to be entered by the user.
For this we need to have a transfer mechanism (a new type of vector) that can
transfer element data arrays. In this way the user can write an element routine
that handles the complicated dependencies of the diffusion coefficient and
stores the result in the vector element by element and is later used by the
diffusion element.

The vector defined per element is of type \texttt{elvector\_t} and can
(currently) store
any number of one-dimensional integer arrays and one-, two- and 
three-dimensional real arrays. The sizes of the arrays can be different for
each element.

The vector needs to be created first, for example:
\begin{listing}{1}
   
  call create_elvector ( mesh, elvector, nint1d=1, nreal3d=2 )

\end{listing}
which creates the vector \texttt{elvector} of type \texttt{elvector\_t} having
one one-dimensional integer array and two three-dimensional real arrays for each
element. Contrary to the creation of sysvectors and vectors during creation of
elvectors no memory is allocated for the data yet. The actual data will be
allocated at the call of the routine \texttt{put\_elvector}:
\begin{listing}{1}

  call put_elvector ( mesh, elvector, elgrp, elem, i1=iarray )

\end{listing}
where the one-dimensional integer array \texttt{iarray} will be stored for
element \texttt{elem} of element group \texttt{elgrp}. The array can be
recovered by
\begin{listing}{1}

  call get_elvector ( mesh, elvector, elgrp, elem, i1=iarray )

\end{listing}
If more than one array of the same data type and dimension has been defined,
the sequence number (if different from 1) needs to be given as well,
for example:
\begin{listing}{1}

  call put_elvector ( mesh, elvector, elgrp, elem, nr=2, r3=array )

\end{listing}
where the three-dimensional real array \texttt{array} will be stored as the
second array. The keyword \texttt{nr} needs to be used for
\texttt{get\_elvector} as well.

If a new array is stored at a position where already an array has been
stored previously, it
depends on the shape of the new array what will happen. If the shape is the
same as the old array, the data is just overwritten. If the shape is different
the old array will be deallocated and the
new one with the new shape will be allocated and stored.

The following statement
\begin{listing}{1}

   call delete ( elvector )

\end{listing}
will delete both all the data arrays and the structure of the array.
The following statement
\begin{listing}{1}

   call delete ( elvector, deletestructure=.false. )

\end{listing}
will delete only the data arrays and leaves the structure intact, as if a fresh
elvector has been created.

Examples of the use of vectors of type \texttt{elvector\_t} are 
\texttt{diffusion2.f90} and 
\texttt{diffusion3.f90} in Sec.\ref{sec:diffusion1}.

\subsection{Subroutines for writing elements that are available in \tfem}
\label{sec:elemsub}

Some subroutines are available in \tfem\ that can help in writing elements. Use
the command \texttt{showmodules} to get an overview of all routines.

At the time of writing routines are available for: getting the coordinates
of the element (in module \texttt{mesh\_m}):

\begin{verbatim}
   get_coordinates
   get_coordinates_sides
   get_coordinates_side
   get_coordinates_geometry
   get_coordinates_object
\end{verbatim}
getting the components and/or positions on element level of
a system vector (in module \texttt{system\_m}):

\begin{verbatim}
   get_sysvector
   get_sysvector_geometry
   get_sysvector_constraint
\end{verbatim}
and vector (in module \texttt{vector\_m}):

\begin{verbatim}
   get_vector
   get_vector_geometry
\end{verbatim}
getting and storing the components of a vector defined per element
for a single element (in module \texttt{elvector\_m}):
\begin{verbatim}
   get_elvector
   put_elvector
\end{verbatim}
Furthermore, there are various routines for shape functions and isoparametric
mapping (in module \texttt{shapefunc\_m}) and for Gauss integration (in module
\texttt{gauss\_m})

\subsection{The structure coefficients}
\label{sec:coefficients}

If the structure coefficients is part of the heading of \texttt{build\_system} (and several other routines),
this structure is passed through to the element subroutine. The structure
coefficients is of type \texttt{coefficients\_t} and must be created with 
the routine \texttt{create\_coefficients}, for example:
\begin{listing}{1}
   
   call create_coefficients ( coefficients, ncoefi=10, ncoefr=20 )

\end{listing}
which creates the structure coefficients having two arrays of length 10 and 20,
respectively.
\begin{verbatim}
   coefficients%i
   coefficients%r
\end{verbatim}
The first is an integer array and the latter is a real array. After filling the
arrays the user can use the data on element level to set model parameters.

Additionally multiple integer and real arrays can be created as follows:
\begin{listing}{1}
   
   call create_coefficients ( coefficients, ncoefi=10, ncoefr=20, &
     ncoefia=[30,40], ncoefra=[15] )

\end{listing}
which creates the structure coefficients having additionally the arrays 
\begin{verbatim}
   coefficients%ia(1)%a
   coefficients%ia(2)%a
   coefficients%ra(1)%a
\end{verbatim}
which are of size 30, 40 and 15, respectively. 

Furthermore, multiple 2D integer and real arrays can be created as follows:
\begin{listing}{1}
   
   call create_coefficients ( coefficients, ncoefi=10, ncoefr=20, &
     ncoefia2=[10,4], ncoefra2=[15,6,20,7] )

\end{listing}
which creates the structure coefficients having additionally the 2D arrays 
\begin{verbatim}
   coefficients%ia2(1)%a
   coefficients%ra2(1)%a
   coefficients%ra2(2)%a
\end{verbatim}
which have shapes (10,4), (15,6) and (20,7), respectively. Note, that the shapes are defined by pairs and that
the number of 2D integer and real arrays created are \texttt{size(ncoefia2)/2} and \texttt{size(ncoefra2)/2}, respectively.

The structure coefficients also contain ``procedure pointers'' for assigning user functions
that need to be used on element level. Available are pointers to scalar (\texttt{func}),
vector (\texttt{vfunc}) and tensor functions (\texttt{tfunc}).
for example:
\begin{verbatim}
   coefficients%func => scalar_func
   coefficients%vfunc => vector_func
   coefficients%tfunc => tensor_func
\end{verbatim}
There are also arrays of these pointers, for example
\begin{verbatim}
   coefficients%func1(1)%p => scalar_func
\end{verbatim}
will assign the function \texttt{scalar\_func} to the first pointer in the
array of pointers \texttt{func1}. The dimension of the arrays is set by
the variable \texttt{MAXFUNCTIONS} in the module \texttt{limits\_m}.
There are more procedure pointers available for special purposes (discussed in the examples).

It is planned to extend the type \texttt{coefficients\_t} to easily provide
different coefficients for different element groups and/or individual elements
in a standardized and general way.
An intermediate mechanism has been implemented for now to
provide different coefficients for different element groups\footnote{Of course,
the same result can be achieved by calling \texttt{build\_system}, \ldots for
different groups separately with different input for the coefficients. However, this
leads to many more statements in your code.}. For example,
assume that there are two elements groups needing different coefficients. Then, define
the following in your main program:
\begin{listing}{1}
  type(coefficients_t) :: mcoefficients(2)

  ! create coefficients

  call create_coefficients ( coefficients(1), ncoefi=100, ncoefr=100 )
  call create_coefficients ( coefficients(2), ncoefi=100, ncoefr=100 )

\end{listing}
Note,
 that the dimension of the array of coefficients \texttt{mcoefficients} should
be equal to the number of element groups.
In the call to the routines \texttt{coefficients} need to be replace with \texttt{mcoefficients}, for example:
\begin{listing}{1}
! stokes velocity/pressure
  call build_system ( mesh, problem, sysmatrix, rhsd, &
    elemsub=stokes_elem, physqrow=[2,3], physqcol=[2,3], &
    mcoefficients=mcoefficients )
\end{listing}
Note, that \texttt{mcoefficients} is a different (optional) heading parameter.
Various routines have been adapted to accept the \texttt{mcoefficients} parameter.
However, only where it makes sense or where the element group number is unique. For example, for curves or surfaces
the group number is not unique and dealing with different coefficients will only be possible after the more
general implementation.

\subsection{The structure oldvectors}
\label{sec:oldvectors}

If the structure oldvectors is part of the heading of \texttt{build\_system},
this structure is passed through to the element subroutine. The structure
oldvectors is of type \texttt{oldvectors\_t} and must be created with 
the routine \texttt{create\_oldvectors}, for example:
\begin{listing}{1}
   
   call create_oldvectors ( oldvectors, nvec=3 )

\end{listing}
This will 
create in the structure oldvectors a component \texttt{v} that
is an array of pointers of length 3. 
Each pointer in the array \texttt{v} can point to a structure of 
type \texttt{vector\_t}. A pointer is encapsulated into a type with a
single component \texttt{p}, hence
\begin{verbatim}
   oldvectors%v(1)%p
   oldvectors%v(2)%p
   oldvectors%v(3)%p
\end{verbatim}
are all pointers to a structure of type \texttt{vector\_t}.

In order to associate a pointer with a target a fortran pointer assignment
statement needs to be performed. For example:
\begin{verbatim}
   oldvectors%v(1)%p => vector
\end{verbatim}
will associate \texttt{oldvectors\%v(1)\%p} with \texttt{vector}. So 
\texttt{oldvectors\%v(1)\%p} can be used as an alias for \texttt{vector}.
In fortran the target (\texttt{vector}) will need to have the \texttt{target}
attribute in the definition:
\begin{verbatim}
   type(vector_t), target :: vector
\end{verbatim}

Similarly, the component \texttt{s} in the structure oldvectors represents an
array of pointers to system vectors (of type \texttt{sysvector\_t}).
If some of
the vectors or system vectors belong to a different problem, the problem
structures can be included in the structure oldvectors:
\begin{listing}{1}
   
   call create_oldvectors ( oldvectors, nsysvec=3, nvec=2, nelvec=1, nprob=2 )

\end{listing}
Here a structure oldvectors is created to store three system vectors
 (in \texttt{oldvectors\%s(1:3)}),
 two vectors
 (in \texttt{oldvectors\%v(1:2)}),
 one vector defined per element
 (in \texttt{oldvectors\%e(1)})
 and two problem structures
 (in \texttt{oldvectors\%p(1:2)}).
Note, that each element of the arrays is of a type that contains a pointer
component \texttt{p}. For example, \texttt{oldvectors\%p(1)\%p} is the
actual (pointer to a) structure of type \texttt{problem\_t}.
Note also that all structures put into oldvectors need to have the
\texttt{target} attribute at its definition.
When dealing with multiple problems it is advised
\begin{itemize}
   \item to number all problems explicitly by setting
      \texttt{input\_probdef\%probnr} to a unique number before the call to
      \texttt{problem\_definition},
   \item to store \emph{all}
        problems in \texttt{oldvectors\%p} in the sequence as you
        have given them numbers. This means that the current problem structure
        is available twice (via the structure \texttt{problem} and
        \texttt{oldvectors\%p}) and not all problem structures might be
        needed on element level, but is much easier and less error prone.
\end{itemize}

Additional storage facilities exist for vectors and
system vectors in the oldvectors type. For example,
\texttt{s1} is a
pointer array where each pointer has a \emph{one-dimensional array of system
vectors} as its target. Similarly \texttt{v1} contains pointers
to one-dimensional arrays of vectors. Also available are the components
\texttt{s2} and \texttt{s3} which consist of pointers to two- and
three-dimensional arrays of system vectors. Example:
\begin{listing}{1}
 
   type(sysvector_t), target :: sol, solv(3), solc(6,3)

   call create_oldvectors ( oldvectors, nsysvec=1, nsysvec1=1, nsysvec2=1 )

   oldvectors%s(1)%p => sol
   oldvectors%s1(1)%p => solv
   oldvectors%s2(1)%p => solc

\end{listing}

Additional storage facilities exist for meshes in the oldvectors type:
component \texttt{m} 
is a one-dimensional array of pointers to the type \texttt{mesh\_t}
For example,
\begin{listing}{1}
   
   call create_oldvectors ( oldvectors, nmesh=1 )

\end{listing}
where now a mesh can be associated with \texttt{oldvectors\%m(1)\%p}.
The mesh needs to have a \texttt{target} attribute at its definition.

Additional storage facilities exist for eltree arrays in the oldvectors type:
component \texttt{ea} 
is a one-dimensional array of pointers to the type \texttt{eltree\_array\_t}
For example,
\begin{listing}{1}
   
   call create_oldvectors ( oldvectors, nelta=1 )

\end{listing}
where now an eltree array can be associated with \texttt{oldvectors\%ea(1)\%p}.
The eltree array needs to have a \texttt{target} attribute at its definition.

\subsection{Subdividing elements for subintegration}
\label{sec:subdivide}

The module eltree provides routines for dividing a simple domain into
smaller subdomains.
The smaller subdomains can also be used to define a new Gauss
integration scheme on (part) of the original simple domain. It uses a levelset function to
separate domains. It is also possible to reconstruct the levelset=0 surface and define
a surface integration scheme on it.

\subsubsection{Eltrees}
\label{sec:eltree}

The module \texttt{eltree\_m} provides routines to divide a
reference element (1D: $[-1,1]$, 2D: $[-1,1]\times[-1,1]$ or
3D: $[-1,1]\times[-1,1]\times[-1,1]$ and similarly for triangle and tetrahedrons)
into subdomains using binary trees
(1D), quadtrees (2D) and octrees (3D). The main purpose is to integrate on
parts of a reference element.
The interface between the parts of a reference element
is defined by a zero levelset function.
If the domain is different from the reference element a mapping function
(\texttt{mapcoor}) needs to be defined.
The smallest subdomains near the interface can be further subdivided using
\texttt{submesh} (see below) into
line elements, triangles or tetrahedrons that are aligned along the zero
levelset.
It is also possible to reconstruct the levelset=0 surface using
\texttt{intmesh} (see below).
Composite Gauss integration schemes can be defined on both the volume and the
zero levelset surface.

Examples are given in Sec.~\ref{sec:exampleeltree}.

\subsubsection{Submesh}
\label{sec:submesh}

The module \texttt{submesh\_m} provides routines to divide a simple
domain into a small number of subdomains consisting
of line elements, triangles or tetrahedrons. 

The main purpose of the module submesh is to divide a simple domain (an
interval, triangle, quadrilateral, tetrahedron or hexahedron) into
subdomains to easily integrate on
parts of the domain. The subdomains are line intervals, triangles or
tetrahedrons.
It uses the subdivision into triangles or tetrahedrons
as described in \cite{Min07}.
A part of the domain is defined by a levelset function and
the subdomains are aligned with the zero levelset.
The number of subdomains is minimal and the interface defined by the
levelset consists of a minimal number of flat surfaces (lines or triangles).
A typical use is in the subdivision of elements for integration in \xfem\ (see
Sec.~\ref{sec:xfem}). If the description of the interface by a flat surface
(or a small number of flat surfaces) is not accurate enough, the module submesh
can be combined with an eltree (see Sec~\ref{sec:exampleeltree}. 

The module is used in some
examples of \texttt{eltree} in Sec.~\ref{sec:exampleeltree}.

\subsubsection{Intmesh}
\label{sec:intmesh}

The module \texttt{intmesh\_m} provides routines to divide a simple
domain into a small number of subdomains consisting
of line elements, triangles or tetrahedrons. It uses the subdivision into triangles or tetrahedrons
as described in \cite{Min07}. The difference with \texttt{submesh} is that
\texttt{intmesh} reconstructs the levelset=0 surface.
The main purpose is to integrate on the levelset=0 surface.
The module is used in some
examples of \texttt{eltree} in Sec.~\ref{sec:exampleeltree}.

\newpage
\section{Programming the main program: special topics}

\tfem\ is run from a main program. In this section some special topics related
to writing main programs will be presented.

\subsection{Essential boundary conditions for nodesets}
\label{sec:essnodeset}

Setting essential conditions for nodesets (see Sec.~\ref{sec:nodesets} for
adding nodesets to the mesh) is similar to curves, \ldots.
Example:
\begin{listing}{1}

  call define_essential ( mesh, input_probdef, nodeset1=2, nodeset2=4, physq=1 )

  ...

  call fill_sysvector ( mesh, problem, sol, &
    physq=1, nodeset1=1, nodeset2=4, func=func, funcnr=3 )

\end{listing}

\subsection{Changing the definition of essential boundary conditions}

In some problems you might want to change
the definition of the essential boundary
conditions. For example, in a time-dependent flow at a particular point in time
you would like to change a certain curve from an essential boundary to a
free boundary. Assume that the following essential boundary conditions are
defined before the change:
\begin{listing}{1}

  call define_essential ( mesh, input_probdef, point=1, physq=2 )
  call define_essential ( mesh, input_probdef, curve1=2, curve2=4, physq=1, &
    degfd=[0,1] )
  call define_essential ( mesh, input_probdef, curve1=1, physq=1 )

\end{listing}
and that we want to change the essential conditions on curve 1 from fully
essential to only essential in $x$ direction and free (natural) in the $y$
direction. This can be achieved by
\begin{listing}{1}
   
  call change_essential ( mesh, input_probdef, bctype='curves', nr=2, &
    degfd=[1,0] )

\end{listing}
Here \texttt{bctype='curves'} means we want to change the definition for
curves and \texttt{nr=2} means we want to change the second definition
for a curve (note
the first call of \texttt{define\_essential} is for a point not a curve).
The parameters \texttt{degfd=[1,0]} means we want to change from fully
essential to only essential in the $x$-direction.
Note, that a call to \texttt{change\_essential} only makes changes to what is
in the heading (no default changes).
Other examples:
\begin{listing}{1}
   
  call change_essential ( mesh, input_probdef, bctype='curves', nr=2, &
    degfd=[0] )

\end{listing}
removes all essential boundary conditions on curve 1,
\begin{listing}{1}
   
  call change_essential ( mesh, input_probdef, bctype='curves', nr=1, &
    degfd=[] )

\end{listing}
makes all degrees (of \texttt{physq=1}) on curve 2 to 4 essential,
\begin{listing}{1}
   
  call change_essential ( mesh, input_probdef, bctype='point', nr=1, &
    physq=3 )

\end{listing}
shifts the essential boundary condition in point 1 from \texttt{physq=2} to
\texttt{physq=3}.

In order to apply the new essential boundary conditions we need to define a new
problem:
\begin{listing}{1}

  call delete( problem )

  call problem_definition ( input_probdef, mesh, problem )
   
\end{listing}
or:
\begin{listing}{1}

  call problem_definition ( input_probdef, mesh, problem2 )
   
\end{listing}
where \texttt{problem2} is new.

The change in essential boundary conditions leads to a new numbering of the
unknowns in \texttt{sysvector}. If solution vectors from previous
time-steps need to
be used after the change in boundary conditions
the old sysvectors need to be transferred to the new
numbering. This can be done as follows:
\begin{listing}{1}

!  system numbering -> nodal numbering

   sysvector%u = sysvector%u( problem%degfdperm(:,2) )

   < change essential and redefine the problem >

!  nodal numbering -> system numbering

   sysvector%u = sysvector%u( problem%degfdperm(:,1) )
   
\end{listing}

\subsection{Boundary conditions}

\label{sec:bounc}

In previous sections we have seen boundary conditions of either
the Dirichlet type or Neumann type. It is also rather common to have a relation
between the normal flux and the solution value (Robin boundary condition),
periodical or  other constrained boundary conditions, boundary conditions in a
rotated coordinate system or
boundary conditions of a special type. In this section we describe how to
handle these boundary condition in \tfem.

\subsubsection{Robin boundary conditions}

As an example we use the Poisson equation \eqref{eq:poisson}. 
Assume the following relation on a part of the boundary ($\Gamma_\text{R}$,
with $\Gamma=\Gamma_D\cup\Gamma_{\text{N}}\cup\Gamma_{\text{R}}$):
\begin{equation}
   -\pderiv{u}{n} = \sigma u + h_\text{R}\text{ on } \Gamma_\text{R}
\end{equation}
where $\sigma\geq0$.
This is called a Robin boundary condition.
We now get the following weak form: Find $u\in S$ such that
\begin{equation}
   (\nabla v,\nabla u) +(v,\sigma u)_{\Gamma_{\text{R}}}+
(v,h_{\text{R}})_{\Gamma_{\text{R}}}+
(v,h_{\text{N}})_{\Gamma_{\text{N}}}=(v,f)
   \label{eq:weakformpoissonRobin}
\end{equation}
for all $v\in V$, where $S$ and $V$ are suitable functional spaces for $u$ and
$v$, respectively. We see that a Robin boundary condition adds a
contribution to both the system matrix and system vector. In \tfem\ there are
basically two ways to achieve this\footnote{A third possibility is to add a
new group of line
or surface elements that connect to the boundary, but this is somewhat
cumbersome.}:
\begin{enumerate}
   \item Using the routine \texttt{add\_boundary\_elements} for building the
system matrix and vector on curves and/or
 surfaces.
   \item Add an object to the mesh (see Sec.~\ref{sec:objects}) using the curve
or surface defining $\Gamma_R$ and use \texttt{build\_system} on this object
to add both the element matrix/vector to the system matrix/vector.
\end{enumerate}

\subsubsection{Periodical or other constrained boundary conditions}

The standard way of treating periodical boundary conditions is using
constraints (see Ch.~\ref{ch:constraints}).

\subsubsection{Boundary conditions in a rotated coordinate system}

The standard way of treating 
boundary conditions in a rotated coordinate system
is using constraints (see Ch.~\ref{ch:constraints}).

\subsubsection{Boundary conditions of a special type}

Sometimes special boundary conditions are required, for example the so-called 
`open boundary condition' used as an outflow boundary condition in the
Navier-Stokes equation. Such boundary conditions can mostly
be handled using objects
(see Sec.~\ref{sec:objects}):  add an object to the mesh and
use \texttt{build\_system} on this object
to add both the element matrix/vector to the system matrix/vector. The
advantage of objects over curves/surfaces is that in addition to the primary
unknowns, also the derivatives in \emph{all} coordinate directions are
available. Also objects can be positioned everywhere in the mesh, not just at
the boundaries.

\subsection{Working with degrees of freedom}

Degrees of freedom are defined in the nodes of the mesh\footnote{Constraints
introduce additional degrees of freedom, the Lagrange multipliers (see
Chapter~\ref{ch:constraints}). The Lagrange multipliers are usually not
defined in the nodes of the mesh and are not considered here.} by the element
degrees of freedom:
\texttt{input\_probdef\%elementdof} for the system vector and 
\texttt{input\_probdef\%vec\_elementdof} for the vectors. If physical
quantities have been defined, \texttt{input\_probdef\%elementdof} does not have
to be specified since it is generated
internally by summing the element degrees of the physical quantities in 
\texttt{input\_probdef\%vec\_elementdof}. See Sec.~\ref{sec:driven_cavity2D}
an example.
   
If no physical quantities have been defined,
the actual storage of the degrees in a node is easy: all degrees one after the
other. This ordering is denoted simply by \textsc{d}.
If physical quantities have been defined, the actual storage for the system
vector of the degrees in
a node is the physical quantities one after the other.
Hence, first all degrees of physical quantity 1 then all degrees of physical
quantity 2, etc. This is denoted by \textsc{pd}, which should be read as two
nested do/for loops:
\begin{verbatim}
  pos = 0
  do physq = 1, nphysq
    do degree = 1, ndegrees(physq)
      pos = pos + 1
    end do
  end do
\end{verbatim}
where \texttt{ndegrees(physq)} are the number of degrees if physical quantity
\texttt{physq} and \texttt{pos} is index/position in the local array of degrees
of freedom. The ordering of the degrees in a node of a vector is of course
still \textsc{d}, since physical quantities only affect the system vector.

The full system vector and the vectors consist of storing the degrees 
node for node\footnote{For the system vector there is an additional
renumbering. For example, essential boundary conditions lead to renumbering, but
there can be more reasons for renumbering.},
thus first all degrees in node 1, then node 2 etc.
For the system vector with physical quantities this is denoted by
the nested loop notation as \textsc{npd},
 i.e.\ the nodes are the `outer loop'. If physical degrees are absent or in the
case of vectors, the ordering of the degrees is \textsc{nd}.

Presenting the element degrees of freedom to the element routines in the
order \textsc{npd} is not optimal. For example, see the element shown in
Fig.~\ref{fig:element}.
\begin{figure}
  \begin{center}
   \includegraphics{element}
  \end{center}
\label{fig:element}
\caption{Degrees in the six-node Taylor-Hood velocity-pressure element.}
\end{figure}
The physical degree 1 is the velocity $(u,v)$ and physical degree 2 is the
pressure $p$. The element vector in the order \textsc{npd} is given by
\[
   (u_1,v_1,p_1,u_2,v_2,u_3,v_3,p_3,u_4,v_4,u_5,v_5,p_5,u_6,v_6)
\]
It is much easier to handle if the physical quantities velocity and pressure are
separated by putting the physical quantities \textsc{p} in the `outer loop'.
For the order \textsc{pnd} this gives
\[
   (u_1,v_1,u_2,v_2,u_3,v_3,u_4,v_4,u_5,v_5,u_6,v_6,p_1,p_3,p_5)
\]
It might also be preferable to separate the $u$ and $v$ components of the
velocities by putting the nodes in the `inner loop'.
For the order \textsc{pdn} this gives
\[
   (u_1,u_2,u_3,u_4,u_5,u_6,v_1,v_2,v_3,v_4,v_5,v_6,p_1,p_3,p_5)
\]
This is the default in \tfem. In fact, physical quantities are always in the
`outer loop'. It is only possible to change the order of the degrees and the
nodes in \tfem\ by the \texttt{order} keyword with the two possibilities
\texttt{'ND'} or \texttt{'DN'}, with an obvious meaning.

To make things even more complicated, in \tfem\ it is possible to define layers
of degrees of freedom (see Sec.~\ref{sec:layers}). 
Denoting 
the layers by \textsc{l}, the actual nodal storage order is \textsc{pld}. Since
the outer loop is \textsc{p}, the physical quantities are separated in
actual storage, even if there exist multiple layers in a node. The actual
storage for multiple nodes is of course \textsc{npld}. In the example above
assuming both layers present in all nodes would lead to the following element
vector for the ordering \textsc{npld}
\begin{multline*}
   (u^1_1,v^1_1,u^2_1,v^2_1,p^1_1,p^2_1,\,
    u^1_2,v^1_2,u^2_2,v^2_2,\,
    u^1_3,v^1_3,u^2_3,v^2_3,p^1_3,p^2_3,\\
    u^1_4,v^1_4,u^2_4,v^2_4,\,
    u^1_5,v^1_5,u^2_5,v^2_5,p^1_5,p^2_5,\,
    u^1_6,v^1_6,u^2_6,v^2_6)
\end{multline*}
where the superscript denotes the layer and the subscript the node. This
ordering is not actually available for the element subroutines. Currently,
there are only two possibilities \textsc{pnld} and \textsc{pldn} as given by
the keyword \texttt{order='ND'} and \texttt{order='DN'}, respectively.
For \textsc{pnld} the example becomes
\begin{multline*}
   (u^1_1,v^1_1,u^2_1,v^2_1,\,
    u^1_2,v^1_2,u^2_2,v^2_2,\,
    u^1_3,v^1_3,u^2_3,v^2_3,,\\
    u^1_4,v^1_4,u^2_4,v^2_4,\,
    u^1_5,v^1_5,u^2_5,v^2_5,\,
    u^1_6,v^1_6,u^2_6,v^2_6,\,\\
    p^1_1,p^2_1,
    p^1_3,p^2_3,
    p^1_5,p^2_5)
\end{multline*}
which separates the physical quantities.
For \textsc{pldn} the example becomes
\begin{multline*}
   (u^1_1,u^1_2,u^1_3,u^1_4,u^1_5,u^1_6,\,
    v^1_1,v^1_2,v^1_3,v^1_4,v^1_5,v^1_6,\\
    u^2_1,u^2_2,u^2_3,u^2_4,u^2_5,u^2_6,\,
    v^2_1,v^2_2,v^2_3,v^2_4,v^2_5,v^2_6,\\
    p^1_1,p^1_3,p^1_5,
    p^2_1,p^2_3,p^2_5)
\end{multline*}
which separates the layers. A third possibility, \textsc{pdln}, separating the 
the $u$ and $v$ components of the velocities leads to
\begin{multline*}
   (u^1_1,u^1_2,u^1_3,u^1_4,u^1_5,u^1_6,\,
    u^2_1,u^2_2,u^2_3,u^2_4,u^2_5,u^2_6,\\
    v^1_1,v^1_2,v^1_3,v^1_4,v^1_5,v^1_6,\,
    v^2_1,v^2_2,v^2_3,v^2_4,v^2_5,v^2_6,\\
    p^1_1,p^1_3,p^1_5,
    p^2_1,p^2_3,p^2_5)
\end{multline*}
would be nice to have, but is not yet available.

If layers are missing from some nodes in the element, the element vector
is reduced accordingly by eliminating the corresponding components from the
previously given element vectors. Similarly rows and columns are eliminated
from the matrix. Therefore the element vectors/matrices change in size from
element to element. This can be rather complicated and other options maybe be
available in a future version of \tfem.

\subsection{Layers}
\label{sec:layers}

In \tfem\ it is possible to define layers of degrees of freedom. For example:
\begin{verbatim}
! problem definition

  call create_input_probdef ( mesh, input_probdef, nvec=3, nphysq=2, &
    numlayers=2 )

  input_probdef%vec_elementdof(1)%a =   &
      reshape ( [2,2,2,2,2,2,    &  ! velocity
                 1,0,1,0,1,0,    &  ! pressure
                 1,1,1,1,1,1, ], &  ! scalar, such as vorticity
                 [6,3] )

  input_probdef%physq = [1,2] ! set elementdof to consist of two physq

  input_probdef%layers = [2,5]  ! nodeset 2 and 5 define the two layers

  ......

  call problem_definition ( input_probdef, mesh, problem )

\end{verbatim}
defines two layers, where layer 1 is present in the nodes given by
nodeset 2 and layer 2 is present in the nodes given by nodeset 5. Each layer
defines a set of nodal degrees as specified by the
element degrees of freedom given by 
\texttt{input\_probdef\%elementdof} for the system vector and
\texttt{input\_probdef\%vec\_elementdof} for the vectors. However, for
\emph{each separate node} there is the choice for a layer being present or
not. So, in the example above we can have four possibilities for a node: no
degrees, only layer 1, only layer 2 or both layer 1 and 2 present. Note, that
a physical quantity is defined by all layers, not just one. 

The usual connections in the \fem\ remain valid, i.e.\ \emph{all} degrees in
a node connect to \emph{all} degrees in \emph{all} nodes of \emph{all}
elements connected to the node. From the element assembly perspective
the full elements connecting all degrees of freedom in an element need to be
assembled.
If the system matrix is build only partially, for example
in a single layer only, this can lead to unexpected/missing
connections between degrees of freedom.
With varying number of layers in the nodes of an element 
the full element vectors/matrices also change in size from
element to element.

There are various possibilities to deal with layers in \tfem, such as:
\begin{itemize}
   \item Essential boundary conditions for a single layer only.
   \item Perform the assembly (\texttt{build\_system}) for a single layer only.
         Use \texttt{derive\_vector} for a single layer only.
         Natural boundary conditions for a single layer only.
         In this way, standard elements can be used.
   \item Define/build constraints (see Sec.~\ref{ch:constraints})
         and connections (see Sec.~\ref{ch:connections}) between layers.
   \item Fill vectors with respect to a single layer only.
   \item Plot a single layer with \texttt{figplot} (Sec~\ref{sec:figplot}).
   \item Sample only a single layer.
\end{itemize}

The primary reason for introducing layers is to deal with material interfaces
in \xfem\ (See Sec.~\ref{sec:xfem}).
The different materials are represented by a single layer and can be build by
standard elements. In the elements at the material interface multiple layers
will be present. The different layers can be coupled at the interface using
either constraints (for weak discontinuities) or connections (for cohesive zone
type interfacing). If no coupling takes place, there is a jump in the field
variable, modeling an internal free surface.

It is also possible to use layers for implementation of the conventional \xfem\
using enrichment. The standard degrees can be in a single layer defined in
all nodes, whereas other layers are only defined in the enriched nodes near the
interface.

For the use of layers we refer to the examples in the section on \xfem\
(Sec.~\ref{sec:xfem}).

\subsection{Limits}
\label{sec:limits}

In the module \texttt{limits\_m} some defaults are set. These are defined as
variables and can be reset at the start of a program. For example, in
\texttt{limits\_m} we have:
   \begin{verbatim}

! limits build in for the coefficients

  integer :: &
    MAXFUNCTIONS = 10  ! default maximum number of pointers to functions
   \end{verbatim}
By setting at the start of program:
   \begin{verbatim}

  MAXFUNCTIONS = 20 
   \end{verbatim}
the number of pointers to functions is increased from 10 to 20.

See the source
file of \texttt{limits\_m} to find out which defaults are set and can be
changed by the user at run time.

\subsection{Excluding partitions with just zeros (zero blocks) from the system matrix}
\label{sec:excludezeroblocks}

Mixed elements usually contain large partitions consisting of just the value zero. Examples are the mixed velocity-pressure Stokes elements, which contain a pressure-pressure zero block (see Sec.~\ref{sec:stokesweak}). Also in implicit time-discretization schemes with multiple physical quantities often not all are coupled to each other, leading to partitions containing the value zero. Although the finite element structure leads to sparse systems, it is assumed that all degrees of freedom are coupled with each other. This results in large parts of the sparse system matrix actually being filled with zeros. Therefore, the storage of the system matrix requires more memory than it could be. Reducing the memory footprint of the matrix is the primary goal of this section, but it could possibly reduce the CPU time for solving the system as well\footnote{For direct solvers the amount of fill-in is important for the time needed to solve the system. Although the matrix storage is smaller, the amount of fill-in might not be or actually might be worse. It is advised to experiment first to see whether excluding zero blocks is useful for the problem at hand. Iterative solvers have not been tested at the time of writing.}.

Zero blocks can be excluded from the system matrix by setting the logical matrix \texttt{input\_probdef\%physqmask}
before the call to \texttt{problem\_definition}. The matrix has a shape of \texttt{(nphysq,nphysq)} and the value of the \texttt{(i,j)} entry of the matrix indicates whether the partition with \texttt{physqrow=i}, \texttt{physqcol=j} should be part of the system matrix. For example, assume \texttt{nphysq=3}. 
Setting \texttt{input\_probdef\%physqmask} as follows\footnote{The matrix entries in \texttt{input\_probdef\%physqmask} are initialized to \texttt{.true.}. Hence, it is not needed to set the whole matrix and only the \texttt{.false.} entries need to be set.}
\begin{verbatim}

  input_probdef%physqmask = reshape ( &
       [ .true.,  .true., . true., &
         .true., .false., .false., &
         .true., .false.,  .true. ], [3,3], order=[2,1] )
         
\end{verbatim}
will indicate that the partitions \texttt{(2,2)}, \texttt{(2,3)} and \texttt{(3,2)} are zero blocks and can be excluded from the system matrix.

The matrix\footnote{The matrix \texttt{input\_probdef\%physqmask} is copied to 
\texttt{problem\%physqmask} in \texttt{problem\_definition}.} 
\texttt{physqmask} affects:
\begin{enumerate}
 \item[a.] the reservation of space in the sparse system matrix,
 \item[b.] the assembling of the sparse system matrix.
\end{enumerate}
   Both a.\ and b.\ are off by default and need to be invoked by the user with an optional
   argument \texttt{usephysqmask=.true.} in \texttt{create\_sysmatrix\_structure} (or 
   \texttt{create\_sysmatrix\_structure\_base}) and
   \texttt{build\_system}, respectively. 
    In this way the user has full control
   when and where \texttt{physqmask} is used. If the assembling
   is partition based (each partition assembled with a separate call), the zero block only needs to be excluded from
   the system matrix (by invoking a.) and is not assembled at all (no call to \texttt{build\_system}):
   \medskip
\begin{verbatim}

! create system matrix

  call create_sysmatrix_structure ( sysmatrix, mesh, problem, &
    usephysqmask=.true. )

\end{verbatim}
   However, if the zero block is part of the element matrix (like the
   pressure-pressure diagonal zero block in the stokes elements), both space reservation and
   the assembling is affected by \texttt{physqmask} and both
   a.\ and b.\ needs to be invoked:
   \medskip
\begin{verbatim}

  input_probdef%physqmask = reshape ( &
       [ .true., .true., &
         .true., .false. ], [2,2], order=[2,1] )
         
  call problem_definition ( input_probdef, mesh, problem )

...

! create system matrix

  call create_sysmatrix_structure ( sysmatrix, mesh, problem, &
    symmetric=.true., usephysqmask=.true. )

...

! build (assemble) matrix and vector from elements 

  call build_system ( mesh, problem, sysmatrix, rhsd, elemsub=stokes_elem, &
    coefficients=coefficients, usephysqmask=.true. )

\end{verbatim}
assuming velocity is \texttt{physq=1} and pressure is \texttt{physq=2}. 

Examples \texttt{stokes1b.f90} and \texttt{stokes11a.f90} are available, showing how optionally the zero pressure-pressure block is excluded from the system matrix.







